\scanpage{100}
degenerate. One can pass \unsure{further|even} to
the henselization (non-strict) of $\underline{\mathcal O}_s$, \unsure{one|there}
finds the situation of a henselian trait
$S$, with residue field $k$ (here $\mathbf Q$ --- or a
finite extension of $\mathbf Q$), of a
curve $U_\eta$ of type $(g,\nu)$ (here, $(1,1)$)
over its generic point $\eta$, one is interested
in the extension $\mathcal E_{U_\eta}$ of $\mathcal E_\eta$ by $\pi$,
where here $\mathcal E_\eta$ still fits into
\[
(30)\qquad 1\longrightarrow T\unsure{\longrightarrow|-}\mathcal E_\eta\longrightarrow\mathcal E_k\unsure{\longrightarrow|-}1
\]
\deleted{We are here \unsure{in characteristic $0$|again}}
\inserted{in addition to}
\[
(31)\qquad \unsure{..|:}\,1-\pi-\mathcal E_{U_\eta}\longrightarrow\mathcal E_\eta\longrightarrow1
\]
Here the outer action of $T$ on
$\pi$ is \inserted{\unsure{?}} surely not trivial (I
conjecture that it is if and only if there
is ``good reduction'') --- thus the group
$\mathcal E_k$ does not itself act outerly
on $\pi$. But every splitting of (30)
\margin{but I have
proved it, in the
context of abelian varieties,
a long
time ago! That must surely
be known on
$H_1$ surely \unsure{!!|!}}
\margin{99}
